% !PRINTED 28
\weberpage{28}
%<<CONT
\begingroup
\setlength{\abovedisplayskip}{5pt plus 2pt minus 1pt}
\setlength{\belowdisplayskip}{5pt plus 2pt minus 1pt}
\setlength{\abovedisplayshortskip}{3pt plus 1pt}
\setlength{\belowdisplayshortskip}{4pt plus 1pt}
\linespread{0.94}\selectfont
gesetzt werden, dass $m\geqq n$ sei und dass $a_0$ und $b_0$ von Null
verschieden sind. Dann ist die Differenz
\begin{equation}
 A-{a_0\over b_0}x^{m-n}B
\tag{2}
\end{equation}
auch eine ganze rationale Function von $x$, deren Grad aber
\gesperrt{kleiner} ist als $m$, da die höchste Potenz in beiden Gliedern
der Differenz denselben Coëfficienten hat und also herausfällt.
Wir setzen diese Differenz:
\begin{equation}
 A'=A'(x)=a'_0x^{m'}+a'_1x^{m'-1}+\cdots,
 \qquad m'<m.
\tag{3}
\end{equation}
Ist nun $m'$ noch $\geqq n$, so können wir in (2) $A'$ an Stelle von
$A$ setzen und dieselbe Schlussweise wiederholen.

So ergiebt sich eine Kette von Gleichungen:
\begin{equation}
\begin{aligned}
 A-{a_0\over b_0}x^{m-n}B&=A',\\
 A'-{a'_0\over b_0}x^{m'-n}B&=A'',\\
 A''-{a''_0\over b_0}x^{m''-n}B&=A''',\\
 &\cdots,
\end{aligned}
\tag{4}
\end{equation}
und diese Kette lässt sich so lange fortsetzen, bis der Grad der
entstandenen Function kleiner als $n$ geworden ist. Da nun in der
Reihe der Functionen $A,A',A''\ldots$ der Grad jeder folgenden
\gesperrt{mindestens} um eine Einheit erniedrigt ist, so besteht die
Kette der Gleichungen (4) \gesperrt{höchstens} aus $m-n+1$ Gliedern;
sie kann aber auch weniger Glieder enthalten, wenn sich gleichzeitig
\gesperrt{mehrere} Potenzen herausheben. Addiren wir die sämmtlichen
Gleichungen (4), bezeichnen die letzte der Functionen $A,A'\ldots$ mit
$C$ und setzen zur Abkürzung
\begin{equation}
 Q={a_0\over b_0}x^{m-n}+{a'_0\over b_0}x^{m'-n}+\cdots,
\tag{5}
\end{equation}
so dass auch $Q$ eine ganze rationale Function von $x$ ist, so folgt
\begin{equation}
 A=QB+C.
\tag{6}
\end{equation}
Die hier geschilderte Operation, durch die aus $A,B$ die Functionen
$Q,C$ gefunden werden, heisst \gesperrt{Division}. $A$ ist der
\gesperrt{Dividendus}, $B$ der \gesperrt{Divisor}, $C$ der
\gesperrt{Rest} und $Q$ der \gesperrt{Quotient}. \gesperrt{Der Grad des
Restes ist immer niedriger als der Grad des Divisors.}
\endgroup
